\documentclass[10pt,openany]{book}
\input{../preamble}
\renewcommand{\thesection}{\arabic{section}}
\renewcommand{\thesubsection}{\thesection.\arabic{subsection}}
\fancyhead[L]{\scriptsize\sffamily はじめての銀河シミュレーション}
\fancyhead[R]{\scriptsize\sffamily スタイルガイド v1.7}
\pretocmd{\section}{\Needspace{9\baselineskip}}{}{}
\begin{document}
\begin{center}
{\Large\sffamily\bfseries\color{navy} Pythonによる\\はじめての銀河シミュレーション}\par\medskip
{\large テスト粒子で学ぶ重力・軌道・銀河相互作用}\par\bigskip
{\Large\sffamily 執筆・実装スタイルガイド}\par\medskip
Version 1.8 / 参考図書整理版 v2.5 対応\\2026年10月8日
\end{center}
\medskip
本ガイドはこれまでの教育方針を引き継ぎ、学生配布資料として第0章を計算環境の準備に限定し、章・節・小節の見出しを動詞句に統一した改訂版です。第0章ではPython環境とJupyter Notebookの利用法だけを扱い、Python文法やNumPy等の一般的な入門は行いません。数値モデルや標準の軌道パラメータを変更する規約ではありません。

\section{本ガイドの目的と教材の基本原則}
各章は、物理の問いを単純なモデルへ置き換え、数式から計算手順を作り、実装して結果を解釈する学習単位とします。
\begin{center}\small\bfseries
物理 → モデル化 → 数式 → 数値計算法 → Python\\→ 数値実験 → 物理的解釈・確認
\end{center}
数値計算法を網羅しません。本書の計算で使う方法だけを本文で導出・実装し、代替法は名前と特徴を短く紹介します。モデルの変更、数値精度の変更、測定条件の変更を混同しないことを全章に共通する原則とします。

\section{本書の位置づけと扱う範囲}
\subsection{読者と到達点}
主な読者は理学部の学部生です。力学、微分積分、ベクトル、Newton重力を学び、Rydenなどの銀河力学ゼミの内容を計算へ移す教材とします。第0章はPython環境の構築とJupyter Notebookの利用法に限定し、Python文法やNumPyの一般的な解説は必要になった章で補います。到達点は、二つの移動するPlummerポテンシャルと質量ゼロのテスト粒子により、潮汐テイル・ブリッジを計算し、条件を変えて比較することです。
\subsection{書名とモデルの対応}
正式書名は表紙・柱・Notebook冒頭・READMEで統一します。「形成」は書名に含めません。固定形状のポテンシャルとテスト粒子による軌道・相互作用が本書の対象であり、宇宙論的な銀河形成計算とは区別します。
\subsection{必須内容にしないもの}
自己重力円盤、大規模N体計算、tree・FMM・GPU・MPI、分布関数からの自己無撞着な初期条件、ガス・星形成・フィードバック、宇宙論的初期条件は扱いません。基本モデルは共面2次元であり、円盤の傾斜角・厚み・鉛直運動も対象外です。三次元化は名称と展望の紹介にとどめます。


\section{正本とソース管理}
\subsection{文章・コード・数値の正本}
文章と数式の正本は\texttt{textbook/chapters/chapXX.tex}、コードの正本はJupytext percent形式の\texttt{code/chapXX/chapXX\_master.py}です。完成版・学生版Notebook、本文のコード抜粋、図、数値表、数値マクロは生成物とします。

本文にPythonコードを手作業でコピーしません。正本のnamed regionから\texttt{generated/chapXX/}へ抽出し、\verb|\inputcode|で読み込みます。\textbf{実行結果の数値にも同じ原則を適用します。} 本文・キャプション・解答の小数値と、実行結果から算出した派生値は、保存データから生成します。図を変えたときだけ本文を手で書き換える運用をしません。

設定値、式の係数、例示のために固定した整数、参考文献の年・ページは生成対象と区別します。丸め桁数は生成側で決め、同じ実測値を別の場所でも同じ定義で表示します。
\subsection{学生版と再生成}
正本の\texttt{BEGIN SOLUTION / END SOLUTION}で解答範囲を囲みます。学生版ではTODOと\texttt{NotImplementedError}へ置き換え、未解答のまま誤った結果へ進ませません。学生版の停止は仕様であり、完成版は全セルを実行して検証します。

全章の更新入口は\texttt{code/rebuild.py}とします。処理順は、正本からNotebook生成、完成版実行、数値・図の更新、本文用生成物の同期、独立テスト、XeLaTeX組版です。生成物だけを修正すると次回失われるため、必ずその正本へ修正を戻します。
\subsection{配布と実行環境}
配布ZIPには、TeX、master、Notebook、標準条件の初期配列、主要な数値データ、図、ビルド手順、検証記録を含めます。フォントファイルは収録しません。実測環境・所要時間・正本のハッシュを記録します。公開先未確定の場合はその旨を明示し、架空の配布URLを書きません。

\section{文体と各章の構成}
\subsection{文体と説明の順序}
本文・注記・キャプションはです・ます調で、一文を短くします。物理的意味を示してから式を置き、式とコードの対応を説明します。必要な専門語は初出で定義し、式や英語名を前提に説明を省きません。
\subsection{章・節・小節の見出し}
本書は学生がPythonで手を動かす実習型教材なので、章・節・小節の見出しは\textbf{原則として動詞句}に統一します。「〜を求める」「〜を計算する」「〜を比べる」「〜を確認する」「〜を区別する」など、学生が何をするのかが具体的に分かる表現を優先します。単に「〜を学ぶ」と付けるのではなく、物理的対象や数値計算法が見える見出しにします。概念・理由を扱う節では「〜の理由を考える」「〜の意味を確認する」のように書きます。

「本章のまとめ」「演習」「演習の解答・着眼点」「参考図書」などの\textbf{定型的な構造見出しは例外}として名詞句のまま残します。図表名、囲み記事名、参考書の正式名称に本則は適用しません。見出しだけを変更する改訂では、数式ラベル・参照先・コード・数値結果を変更しません。

\subsection{固定する構造}
第1--7章は、目標、問い、モデルと式、数値計算法、計算手順、Python実装、数値実験と解釈、まとめ・演習を基本とします。既習の導出は必要以上に繰り返さず、参照先を示します。第0章は準備章として別扱いにし、環境構築とNotebook操作だけを記述します。

第1--7章冒頭の「この章で学ぶこと」は\texttt{goalbox}に統一し、到達目標を3--5項目で書きます。第1--7章末は「本章のまとめ」「演習」で終えます。第0章にはgoalbox・まとめ・演習を置きません。演習の解答・着眼点は第1--7章分を巻末へ集約し、参考図書も巻末へまとめます。


\section{数式と記法}
\subsection{ベクトル、粒子、銀河中心}
ベクトルは太字、大きさは細字です。銀河中心を$\boldsymbol R_k,\boldsymbol V_k$、テスト粒子を$\boldsymbol r_i,\boldsymbol v_i$とし、$\boldsymbol D=\boldsymbol R_2-\boldsymbol R_1$、$\boldsymbol U=\boldsymbol V_2-\boldsymbol V_1$を相対運動に使います。$M_k,a_k$はPlummer場、$b$は中心間有効力の長さです。
\subsection{衝突しやすい記号を分ける}
質量比を$q=M_2/M_1$に限定し、円盤の切断比は$x_{\max}=R_{\max}/R_d$とします。回転符号は$s_{\rm rot}=\pm1$で、コードの\texttt{spin}に対応します。局所位置$\boldsymbol s_{k,i}$と出身中心からの距離$s_i$は、太字と大きさの関係を明示します。積分のダミー半径は$r'$または$R'$を使います。記号は各章で初出時に定義します。章をまたぐ記号の意味を変更しないようにし、必要なら巻末索引等で一覧化します。第0章には記号一覧を置きません。
\subsection{時間積分}
KDKは$\boldsymbol r^n,\boldsymbol v^n,\boldsymbol v^{n+1/2}$を用い、半整数の速度は中間の更新値です。履歴と保存量には最後のkickを終えた整数時刻の速度を使います。RK4は状態$\boldsymbol y=(\boldsymbol r,\boldsymbol v)$と傾き$\boldsymbol k_1,\ldots,\boldsymbol k_4$を使い、傾きに$\Delta t$を含めません。
\subsection{単位と番号}
内部単位は距離kpc、時間Gyr、質量$\Msun$、速度kpc\,Gyr$^{-1}$です。表示でkm\,s$^{-1}$へ換算する場合は明示します。使用する換算係数と$G$は第1章で導入し、以後の章で同じ値を使います。単位を変える場合は基準量と変換を説明します。

図・表・数式は章番号付きにします。第0章で表を用いる場合も0.1のように章番号を付けます。TeXの太陽質量は\verb|\Msun|で統一します。ラベルは\texttt{fig:ch...}、\texttt{eq:ch...}、\texttt{tab:ch...}など内容を表す名前とし、図番号を本文へ手で書きません。

\section{Pythonコード規約}
\subsection{配列と抽象化}
位置・速度は$(N,2)$、状態は$(N,4)$、履歴は$(N_{\rm save},N,4)$です。中心を行に含める場合は、その先頭行と粒子番号の対応を説明します。\texttt{axis}、スライス、\texttt{[:, None]}は配列の形と物理量を対応させます。

新しい式は最初に成分で示し、必要になってからループ・配列演算へ進みます。学習済みの処理を短い関数へまとめ、本文でクラス設計へ広げません。加速度関数と積分器を分離し、モデルだけを交換できる構造を保ちます。
\subsection{書式と依存}
名前は\texttt{snake\_case}、コメントは日本語を基本にします。型ヒントは必須ではありません。単純な代入・積和・関数呼び出しを不必要に改行しません。長い完全実装はNotebookへ置き、本文の抜粋は原則10--30行程度とします。

計算の中心はNumPy・Matplotlibと標準ライブラリです。新たなライブラリが必要なら理由を記します。シェルの実行コマンドはPythonセルと区別し、共通の\verb|\inputterminal|で表示します。Pythonの\texttt{verbatim}手書き掲載は行いません。


\subsection{主要数値計算法と代替法}
\begin{center}\small
\begin{tabularx}{\linewidth}{lXX}\toprule
用途 & 本文で説明する方法 & 短い補足に留める例\\\midrule
時間更新の入口 & Euler法 & ---\\
重力軌道 & leapfrog（KDK） & 高次シンプレクティック法\\
一般ODE・速度依存力 & 古典的RK4 & RK2、適応刻みRK、Bulirsch--Stoer\\
定積分 & 台形則 & Simpson則、Gauss求積\\
求根 & 二分法 & Newton法、割線法、Brent法\\
補間 & 線形補間 & 多項式、三次スプライン、PCHIP\\
標本化 & CDFの逆変換 & 棄却法、Monte Carlo法の一般論\\\bottomrule
\end{tabularx}\end{center}
KDKは加速度を再利用し、$N_{\rm step}$ステップに$N_{\rm step}+1$回の新しい加速度評価を使います。RK4は各ステップ4回の右辺評価です。同じ刻みでの比較と同じ力評価回数での比較を分け、「高次だから常に有利」「低コストだから常に高精度」とは説明しません。

\section{銀河モデルと数値実験の共通規約}
\subsection{Plummer場とテスト粒子}
銀河の標準場は$\Phi=-GM/\sqrt{r^2+a^2}$とします。星はこの場を受けますが、星同士と星から銀河中心への重力は含めません。粒子数を増やしても背景質量は増えません。指数円盤の粒子数面密度とPlummerの質量密度は独立の量です。初期外縁$R_{\max}$を反射境界や除去境界には使いません。
\subsection{中心間力}
中心の有効対ポテンシャルは$W_c=-GM_1M_2/\sqrt{D^2+b^2}$です。共通の$b$から二つの力を求め、作用・反作用を保ちます。これは広がった二つのPlummer密度分布の厳密な相互エネルギーではありません。標準の$b=a_1=a_2$では、中心への相手場と星への相手場が対応し、潮汐加速度差を同じ場から書けます。

$b$を変えるのは物理モデルの変更であって、刻みを細かくする数値収束とは別です。同じ距離での点質量との力の比は、潮汐力や最終形態の比ではありません。標準の$D_{\rm p}<R_{\max}$は強い重なりを持つ遭遇なので、遠方の弱い潮汐展開を全円盤へ適用しません。
\subsection{動力学的摩擦}
第7章で、固定Plummer主銀河・質量比1:10の衛星へ切り替えます。衛星中心にChandrasekhar加速度を与え、既習のRK4を使います。摩擦を質量ゼロの星へ直接加えません。背景応答を直接解かないこと、局所Maxwell近似、固定Coulomb対数、衛星の質量損失を省くことを明示します。

速度依存力がleapfrog系の方法では原理的に一切扱えない、とは説明しません。本書の標準的な陽的KDKをそのまま適用しない、という範囲を明確にします。


\section{図と可視化のスタイル}
\subsection{数値図と概念図}
数値結果の図はmasterと保存データから生成します。概念・幾何・操作順序はTikZを使い、\texttt{figures/tikz/}へ1図1ファイルで保存します。台形則など元の数値図へTikZ注記を重ねる場合も、元図を残します。

本文には図への導入文と日本語キャプションを付けます。図中は短い英語・本文と同じ数式記号を使います。概念図の説明用配置を、標準計算のスナップショットや測定値と見せません。
\subsection{配色・縮尺・フォント}
白背景とOkabe--Ito系の配色を基本とします。主銀河は青（\texttt{0072B2}）、伴銀河は朱色（\texttt{D55E00}）、中心は黒、比較線は灰色または線種で区別します。色だけで意味を伝えず、マーカーと線種も併用します。軌道や円盤の縦横比は等しくし、比較図は原則同じ軸範囲にします。

TikZの基本文字は9pt、短い補助ラベルは8--8.5ptです。\verb|\resizebox|による図全体の拡大縮小で文字サイズを変えず、座標・余白を調整して本文幅に収めます。スタイルは\texttt{preamble.tex}に一元定義し、図単体の読み込み順で未定義にならないようにします。
\subsection{アニメーションと出力}
アニメーションを補助として使って構いませんが、主要結果は静止画と数値でも読めるようにします。再生間隔・履歴保存間隔・時間積分の刻みを区別します。出身粒子の色とラベルは途中で変更しません。

\section{囲み記事・演習・Notebook}
\subsection{囲み記事と章の構造要素}
囲み記事は\texttt{failbox}（よくある失敗）と\texttt{notebox}（もう少し詳しく）の2種類です。前者は実際の誤解や実装の落とし穴、後者は本文の主線に必須でない補足を扱います。代替手法は名前、基本的考え方、特徴または不採用理由、詳しい参照先に限定します。

\texttt{goalbox}は囲み記事ではなく、章冒頭に1回置く構造要素です。3--5個の目標を統一書式で示します。\texttt{pointbox}を別用途で増やしません。短い確認は本文中の「確認：」で示します。
\subsection{演習と解答}
確認・応用・発展の3段階を第1--7章で使います。確認は本文の理解、応用は既習の手法による条件変更、発展は新しい比較や物理への拡張です。第0章には演習を置きません。解答は巻末へ集約し、数値と着眼点に留めます。完全コードは完成版Notebookで確認できるようにします。
\subsection{Notebook}
本文の長い理論説明を丸ごと再掲しません。問題設定、式の対応、実装、図、短い問いを並べます。全章のNotebookを独立に上から実行できる構成にし、前章の隠れたカーネル状態に依存させません。


\section{確認・検証と物理的解釈}
各章で少なくとも一つの独立な確認を持ちます。円軌道なら同時刻の解析位置、二体系なら作用・反作用と重心、質量積分なら解析式、CDFなら表の外の独立な比較値を使います。解析解がない軌道は、同じ初期配列・同じ評価時刻で自己収束を確認します。

保存するはずの量だけを保存量と呼びます。初期値が0ならその値で相対誤差を作りません。テスト粒子の表示上の重みを中心二体系のエネルギーへ足して全系の保存量とすることもしません。摩擦ではエネルギー保存でなく仕事と変化の収支を確かめ、後処理の台形則の誤差を分けます。

数値モデル、数値精度、有限標本、測定閾値・時刻を別々に記録します。閾値を越えた粒子割合を非束縛率や質量損失率とは呼びません。非単調な結果も残し、比較だけで因果関係を断定しません。共面条件での結果を任意の傾斜へ一般化しません。

実行時間は内部の検証記録として測定しても構いませんが、学生向け本文の第0章には著者環境の待ち時間表を置きません。性能比較を教材で扱う場合だけ、環境と計時範囲を明示します。

\section{執筆時チェックリスト}
\begin{itemize}
\item 章・節・小節の見出しは原則として動詞句になり、学生が行う計算や確認が具体的に分かりますか。構造見出しは定型名を保っていますか。
\item 第1--7章では問いとモデルの仮定が先に示され、目標が3--5項目になっていますか。第0章は準備だけに限定されていますか。
\item 本文の記号・内部単位・配列の軸がコードと一致していますか。
\item コード抜粋と数値・派生値は正本から再生成できますか。
\item 初期値と評価時刻を固定した比較になり、誤差・標本・測定条件が分離されていますか。
\item 図・表の番号と参照は自動で解決し、第0章も0.xになっていますか。
\item 図中文字は原寸で読め、TikZの矢印・ラベルが重なっていませんか。
\item 第1--7章末が「まとめ→演習」で終わり、巻末の解答・参考図書へ正しく集約されていますか。第0章に不要なまとめ・演習が残っていませんか。
\item 完成版は新しいカーネルで全セル実行し、独立テストに通っていますか。
\item PDFを実データの図表で組版し、ページ画像で確認しましたか。
\end{itemize}


\subsection{参考図書の配置と選び方}
参考図書は各章末へ重複して置かず、巻末の「参考図書」に集約します。銀河力学・天体力学はRydenとBinney \& Tremaineを基本とします。数値計算法は、まず情報源にある和書を用い、浮動小数点数・補間・台形則・求根・Runge--Kutta法は杉浦、逆補間・乱数・逆変換サンプリングは伊理・藤野、leapfrogと天文学の粒子計算は『シミュレーション天文学』を主に参照します。Python 3とNumPy・SciPyを用いたEuler法・古典的RK4・補間・数値積分の実装例をさらに確認する補助参考書として、幸谷『Python数値計算プログラミング』を挙げます。

本書の範囲で和書に同等の説明がある方法について、Landauの\textit{Computational Physics}、ButcherやHairerの常微分方程式専門書を重複して基本参考図書へ挙げません。和書では不足する内容を本文へ追加するときに限り、その内容を確認したうえで洋書を追加します。Python・NumPy・SciPy・JupyterのAPI仕様は、巻末の「補足文献・公式ドキュメント」へ分けます。

\section{Version 1.8で固定する規約}
書名・副題、0--7章構成、KDKとRK4の役割、実行結果の生成値化、第1--7章の章冒頭goalbox、2種類の囲み記事、第1--7章末を「まとめ・演習」に統一し、演習解答と参考図書を巻末へ集約する構成を固定します。第0章は環境構築とJupyter Notebookの利用法だけを扱う準備章とします。未使用の\texttt{pointbox}や部構成の書式は削除します。

基本軌道のパラメータは全章改訂前を維持します。改訂のために実験系列を選別したり、非単調な結果を消したりしません。配布先の公開URLは未確定事項として残し、刊行時に一度だけ確定・反映します。
\end{document}
